The expansion of autonomous exploration on Mars has unlocked unprecedented scientific frontiers, yet it requires navigating a landscape characterized by frequent landslides that drive geomorphological change by mass movement of rock and debris. This presents critical mechanical hazards to robotic assets. Mitigating these risks through accurate landslide prediction is therefore essential to safeguarding mission longevity and preventing catastrophic loss. While landslide prediction from multispectral images using traditional methods yields high accuracy, its generalizability is limited. In this work, we study the performance of self-supervised learning-based (SSL) image models, specifically trained as joint-embedding predictive architectures (JEPA) and develop a novel physics-aware image encoder, showing additional performance gains while maintaining the generalizability provided by SSL models. Experiments reveal that our physics-aware masking strategy captures long-range dependencies between spectral bands. We show that by learning cross-spectral latent correlations rather than raw pixel reconstructions, our model is significantly more robust to the sensor noise and varying lighting conditions commonly encountered in multispectral remote sensing data. We train JEPA encoders employing different masking strategies, and benchmark their downstream landslide segmentation performance on the MMLSv2 dataset. This work directly improves autonomous safety for planetary exploration and establishes a scalable framework for self-supervised learning in high-dimensional remote sensing and terrestrial disaster response.
